\documentclass[12pt, a4paper]{article}

\usepackage[margin=1in]{geometry}
\usepackage{amsmath, amssymb}
\usepackage{graphicx}
\usepackage{booktabs}
\usepackage{hyperref}
\usepackage{xcolor}
\usepackage{listings}
\usepackage{parskip}
\usepackage{titlesec}
\usepackage{fancyhdr}
\usepackage{caption}
\usepackage{float}

\hypersetup{
    colorlinks=true,
    linkcolor=blue,
    urlcolor=blue,
    citecolor=blue
}

\lstset{
    basicstyle=\ttfamily\small,
    backgroundcolor=\color{gray!10},
    frame=single,
    breaklines=true,
    columns=fullflexible
}

\pagestyle{fancy}
\fancyhf{}
\rhead{\small Speculative Decoding -- Prashant Gautam}
\lhead{\small IEOR, IIT Bombay}
\rfoot{\thepage}

\titleformat{\section}{\large\bfseries}{}{0em}{}[\titlerule]
\titleformat{\subsection}{\normalsize\bfseries}{}{0em}{}

\begin{document}

% ---------------------------------------------------------------
% Title Page
% ---------------------------------------------------------------

\begin{titlepage}
    \centering
    \vspace*{2cm}

    {\LARGE \textbf{SpecDecode: Building a Speculative Decoding}}\\[0.4em]
    {\LARGE \textbf{Inference Accelerator from Scratch}}\\[2em]

    {\large Prashant Gautam}\\[0.4em]
    {\normalsize Department of Industrial Engineering and Operations Research}\\
    {\normalsize Indian Institute of Technology Bombay}\\[0.4em]
    {\normalsize Roll No. 24B4526}\\[1em]

    {\normalsize
        \href{https://github.com/PrashantGautam-alt/specdecode}
             {github.com/PrashantGautam-alt/specdecode}
    }\\[3em]

    {\normalsize June 2026}

    \vfill

    \begin{abstract}
        \noindent
        Large language models generate text one token at a time, with each token
        requiring a full forward pass of the target model. This sequential dependency
        is the fundamental throughput bottleneck at inference time. This project
        implements two families of acceleration from scratch on Llama 3.1/3.2 models
        and measures every result on an RTX A5000 (24\,GB): classical speculative
        decoding with a separate draft model, and Medusa, where small heads on the
        frozen backbone predict several future tokens at once. Speculative decoding
        with $K{=}4$ reaches 1.17x. The Medusa path is driven through a sequence of
        structural optimizations -- candidate \emph{tree} attention, a \emph{fused}
        single-pass decoder, and \emph{typical} acceptance -- reaching a provably
        \textbf{lossless 1.24x} (verified by a self-consistency test, not a fragile
        bit-match, because fp16 ties make exact matching meaningless) and a coherent
        \textbf{1.47x} under typical acceptance. The central contribution is a
        \textbf{calibrated candidate tree}: per-head prediction accuracies are
        measured offline and a fixed node budget is allocated by Prim's-style greedy
        selection over path probability. The optimizer \emph{automatically} drops the
        weakest head (assigning it zero nodes -- an emergent re-derivation of the
        optimal-$K$ result), lifting the lossless speedup to \textbf{1.34--1.56x} at
        the same compute budget. A FastAPI + WebSocket service streams tokens to a
        browser, colour-coded live by whether each token was a confirmed speculative
        guess or a backbone correction.
    \end{abstract}
\end{titlepage}

% ---------------------------------------------------------------
\section{Motivation}
% ---------------------------------------------------------------

The standard autoregressive loop is simple: run the model, take the last logit,
sample a token, append it, repeat. At 38 tokens/sec on a Llama 3.1 8B model,
generating a 500-token response takes about 13 seconds. This is slow enough to
matter for user-facing applications.

The key observation in Leviathan et al. \cite{leviathan2023} is that the bottleneck
is not compute but \textit{sequential dependency}. The transformer can process
multiple tokens in parallel during training (each position attends to all previous
positions in one pass), but during generation we feed one token at a time because
each token depends on the previous one. Speculative decoding breaks this dependency
partially by letting a cheap draft model propose several tokens at once, then
verifying all of them in a single target model pass.

% ---------------------------------------------------------------
\section{Background}
% ---------------------------------------------------------------

\subsection{Autoregressive Generation}

Given prompt tokens $x_1, \ldots, x_n$, the model generates token $x_{n+1}$ by
computing:
\[
    x_{n+1} \sim p_\theta(\cdot \mid x_1, \ldots, x_n)
\]
Each step requires a full forward pass. With the KV cache, keys and values for
positions $1 \ldots n$ are stored from the previous step, so only the new token
needs to be processed. This makes each step $O(1)$ in compute but the sequential
nature remains -- step $t+1$ cannot begin until step $t$ produces a token.

\subsection{The KV Cache}

Without caching, generating $T$ tokens on a sequence of length $n$ costs
$O(T \cdot n)$ forward passes in terms of attention operations. With the KV cache,
keys and values for all seen positions are stored and reused. Each new step only
computes $Q$ for the newest token, reduces attention operations to $O(T)$. The
tradeoff is memory: for Llama 3.1 8B with sequence length 2048, the KV cache
occupies roughly $2 \times 32 \times 2048 \times 128 \times 2 \approx 0.5$ GB
per batch element in float16.

% ---------------------------------------------------------------
\section{Speculative Decoding}
% ---------------------------------------------------------------

\subsection{Algorithm}

Speculative decoding \cite{leviathan2023} uses two models: a small fast draft model
$q$ and a large accurate target model $p$. Each round:

\begin{enumerate}
    \item \textbf{Draft:} Run the draft model autoregressively for $K$ steps,
          collecting tokens $\tilde{x}_1, \ldots, \tilde{x}_K$ and their
          distributions $q(\cdot \mid \cdot)$.
    \item \textbf{Verify:} Feed all $K$ draft tokens to the target model in one
          forward pass. The transformer processes all positions in parallel,
          yielding $K$ distributions $p(\cdot \mid \cdot)$ simultaneously.
    \item \textbf{Accept:} For each draft token $\tilde{x}_i$ from left to right,
          accept with probability $\min(1, p(\tilde{x}_i) / q(\tilde{x}_i))$.
          Stop at the first rejection.
    \item \textbf{Correct:} On rejection, sample the correction token from
          $\text{normalize}(\max(0, p - q))$.
\end{enumerate}

\subsection{Rejection Sampling Correctness}

The acceptance criterion $\min(1, p(x)/q(x))$ guarantees that the marginal
distribution of each accepted token equals $p$ exactly. This is not an
approximation -- speculative decoding is lossless by construction.

Proof sketch: let $\alpha = \min(1, p(x)/q(x))$. The probability of outputting
token $x$ is:
\[
    q(x) \cdot \alpha + (1 - A) \cdot r(x)
\]
where $A = \sum_x q(x) \min(1, p(x)/q(x))$ is the overall acceptance probability
and $r(x) = \max(0, p(x) - q(x)) / (1 - A)$ is the residual distribution.
Expanding and simplifying yields $p(x)$ for all $x$.

Empirically verified with a 100K trial test: maximum deviation 0.27\% from target
distribution $p$ -- within the expected $O(1/\sqrt{N})$ Monte Carlo error.

\subsection{Cost Model}

Each round costs $K$ draft passes plus 1 target pass. If a draft pass costs
fraction $c$ of a target pass:
\[
    \text{speedup} \approx \frac{M}{c \cdot K + 1}
\]
where $M$ is the average tokens accepted per round. Empirically, $c \approx 0.40$
on the A5000 -- each 1B draft pass costs about 40\% of an 8B target pass.
This is higher than the parameter ratio ($1/8 = 0.125$) because single-token
forward passes are dominated by fixed per-call overhead, not by FLOPs.

At K=4 with $M = 3.05$: predicted speedup $= 3.05 / (0.40 \times 4 + 1) = 1.17$x.
This matched the measured value exactly.

% ---------------------------------------------------------------
\section{K Sweep Results}
% ---------------------------------------------------------------

\begin{table}[H]
\centering
\caption{K sweep. Draft: Llama 3.2 1B Instruct. Target: Llama 3.1 8B Instruct.
         100 tokens, 3 runs per K. Naive 8B baseline: 38.0 tok/s.}
\begin{tabular}{@{}ccccc@{}}
\toprule
K & tok/s & Speedup & Avg tokens/round & Predicted speedup \\
\midrule
1 & 27.8 & 0.73x & 1.00 & 0.71x \\
2 & 38.4 & 1.01x & 1.83 & 1.02x \\
\textbf{4} & \textbf{44.4} & \textbf{1.17x} & \textbf{3.05} & \textbf{1.17x} \\
6 & 39.0 & 1.03x & 3.65 & 1.07x \\
8 & 34.0 & 0.89x & 3.82 & 0.95x \\
\bottomrule
\end{tabular}
\end{table}

K=4 is optimal because draft cost grows linearly and unconditionally (+0.40 per
draft token regardless of acceptance), while acceptance gains are sublinear and
conditional (one wrong token wastes all remaining draft passes that round). The
cost model captures both effects and predicts the peak correctly.

A key finding: using the base 1B model as draft against the instruct 8B target
peaked at only 0.97x. Switching to the instruct 1B draft raised acceptance from
2.45 to 3.05 at K=4. Base and instruct models operate in different ``dialects'' --
base continues text, instruct responds like an assistant. Matching the training
style of draft and target is as important as the model size ratio.

% ---------------------------------------------------------------
\section{Medusa}
% ---------------------------------------------------------------

\subsection{Architecture}

Medusa \cite{cai2024} eliminates the separate draft model entirely by attaching
small MLP heads directly to the frozen backbone. Each head reads the final hidden
state $h \in \mathbb{R}^{d}$ (where $d = 4096$ for Llama 8B) and predicts a token
further in the future:
\[
    \text{head}_k(h) = W_2^{(k)}\bigl(\text{SiLU}(W_1^{(k)} h) + h\bigr)
\]
Head $k$ predicts the token at position $t + k + 1$. The residual $+h$ is the key
structural choice: at initialization ($W_1 = 0$), the head reduces exactly to the
LM head, giving a warm start rather than random noise.

\subsection{Initialization Trick}

Setting $W_1 = 0$ and $W_2 = \text{LM head copy}$ at initialization ensures every
head starts as an exact clone of the LM head. Verification: $\text{SiLU}(W_1 h) =
\text{SiLU}(0) = 0$, so the bracket reduces to $h$, and $W_2 h = \text{LM head}(h)$.
Confirmed at runtime with \texttt{torch.allclose} to tolerance $10^{-3}$ (float16
precision limit).

\subsection{Training}

The backbone is frozen; only the $K = 4$ heads train. For head $k$ with shift
$s = k + 1$, labels are $\text{input\_ids}[:, s:]$ and logits are
$\text{head\_logits}[k][:, :-s, :]$. The total loss is:
\[
    \mathcal{L} = \sum_{k=0}^{K-1} 0.8^k \cdot \mathcal{L}_k
\]
The $0.8^k$ weighting down-weights far heads because (1) their loss is noisier and
(2) acceptance is sequential -- a wrong head-0 prediction ends the round regardless
of how accurate heads 1--3 are. Trained on UltraChat 25k examples, 2 epochs,
backbone on \texttt{cuda:0} in float16, heads on \texttt{cuda:1} in float32
(float16 heads produce NaN due to softmax overflow at vocab size 128,256).
8-bit Adam (bitsandbytes) used for optimizer memory efficiency.

\subsection{Greedy Decoding Results}

\begin{table}[H]
\centering
\caption{Three-way benchmark. Hardware: 2x RTX A5000. 100 tokens, 3 runs.}
\begin{tabular}{@{}llcc@{}}
\toprule
Config & Notes & tok/s & Speedup \\
\midrule
Naive 8B baseline & & 37.9 & 1.00x \\
SpecDecode K=4 & 1B instruct draft & 44.4 & 1.17x \\
Medusa greedy (float32 heads) & epoch1 checkpoint & 24.6 & 0.65x \\
Medusa greedy (float16 heads) & epoch1 checkpoint & 26.4 & 0.69x \\
\bottomrule
\end{tabular}
\end{table}

This first Medusa implementation is \emph{slower} than the naive baseline. The next
section explains why (it is structural, not a tuning problem), and the sections after
that remove the bottleneck and turn Medusa into the project's fastest configuration.

% ---------------------------------------------------------------
\section{Why Greedy Medusa is Slow: The 3-Pass Problem}
% ---------------------------------------------------------------

Medusa greedy with 2.22 tokens/round acceptance should be faster than naive, but
it runs at 0.69x. The reason is structural: our greedy implementation does 3
backbone passes per round:

\begin{enumerate}
    \item \textbf{PROPOSE:} Feed last token, get hidden state $h$ for the heads,
          get backbone prediction at position 0.
    \item \textbf{VERIFY:} Feed $K-1$ candidates to get backbone predictions at
          positions 1 through $K-1$.
    \item \textbf{CACHE UPDATE:} Re-feed accepted tokens into a clean cache snapshot
          to advance state correctly.
\end{enumerate}

The cache update is forced by a \texttt{DynamicCache} mutation bug in transformers
4.56.2: the KV cache is modified in place during the verify pass, contaminating it
with rejected tokens. The solution -- snapshotting the cache before verify -- is
correct but adds the third pass.

Expected speedup from the cost model: $2.22 / 3 \approx 0.74$x. Measured 0.69x,
with the gap explained by cross-GPU PCIe latency ($h$ crosses \texttt{cuda:0}
$\to$ \texttt{cuda:1} every round).

% ---------------------------------------------------------------
\section{Tree Attention}
% ---------------------------------------------------------------

\subsection{Motivation}

The 3-pass bottleneck is not caused by poor acceptance -- 2.22 tokens/round is
reasonable. It is caused by verification structure. Even with perfect acceptance
(K=4 every round), 3 passes gives at most $4/3 \approx 1.33$x. Tree attention
collapses all 3 passes into approximately 1.

\subsection{Construction}

Instead of one candidate chain $[c_0, c_1, c_2, c_3]$, keep the top-2 from each
head. This yields a binary tree of depth 4: $2 + 4 + 8 + 16 = 30$ nodes,
$2^4 = 16$ candidate paths.

All 30 nodes are flattened into a single sequence and fed to the backbone in one
pass with a custom attention mask. The mask enforces: node $i$ can attend to node
$j$ if and only if $j$ is an ancestor of $i$ in the tree (or $j = i$). Siblings --
nodes on different branches -- are blocked from attending to each other because they
represent mutually exclusive futures. Allowing cross-branch attention would
contaminate the backbone's logit for node $c$ (child of $a$) with information from
branch $b$, producing a logit that does not correspond to any real path.

The full attention mask has shape $[1, N, L + N]$ where $L$ is the context length
and $N = 30$. The left $L$ columns are all ones (every candidate can see the full
context). The right $N \times N$ block is the tree mask built from parent indices.

\subsection{Result}

Tree attention with strict-greedy acceptance is lossless and removes the verification
redundancy, but on its own it still pays for a separate PROPOSE pass. The decisive step is
fusing PROPOSE into VERIFY (next section), after which the tree decoder becomes the
project's fastest lossless configuration.

% ---------------------------------------------------------------
\section{The Fused Single-Pass Decoder}
% ---------------------------------------------------------------

The standalone PROPOSE pass existed only to produce a fresh hidden state $h$ for the heads.
But the VERIFY pass already computes a hidden state at \emph{every} tree node, and these were
being discarded. The fused decoder keeps them: the hidden state at this round's last accepted
node is exactly what the heads need to build the \emph{next} round's tree. This folds PROPOSE
into VERIFY and absorbs the cache update, giving \textbf{one backbone pass per round}.

Fusion did not double the speed (1.19x $\to$ 1.24x). The fused pass prepends the already-known
pending token, which consumes head~0, so the fused tree has depth $K-1$ instead of $K$ -- one
fewer level of guesses. On a memory-bound model a backbone pass is worth roughly one tree
level, so removing a pass while losing a level nearly cancels. The real unlock is therefore
\emph{more heads}, so a fused tree can be deep and still cost one pass.

% ---------------------------------------------------------------
\section{Correctness Under fp16: Ties and Self-Consistency}
% ---------------------------------------------------------------

The fused output disagreed with HuggingFace greedy on a single token (`` describes'' vs.\
`` explains''). This was first chased as a cache bug, but a fresh cacheless recompute matched
the cached result -- the cache was correct. The true cause was a near-\emph{tie}: the two
tokens had logits $19.640$ and $19.625$. The argmax of a tie is computed in different
floating-point reduction orders by a \emph{prefill} pass (parallel over many positions, as the
fused decoder runs) versus a \emph{decode} pass (one token at a time, as HuggingFace greedy
runs); in float16 the last-digit rounding differs and flips the tie.

The consequence is methodological: the fused decoder produces a \emph{valid} greedy decode, but
not one that is bit-identical to an incremental reference. Correctness is therefore tested by
\textbf{self-consistency}: take the fused output, run one clean forward pass over it, and
confirm every committed token is the model's argmax or tied for it within a tiny logit gap.
Result: \textbf{101 / 101 tokens argmax-or-tied}, largest gap $0.0000$ -- the output is provably
a valid greedy decode. Lossless, at \textbf{1.24x}.

\textit{A safety corollary used throughout:} the heads only \emph{propose}; the backbone
\emph{verifies and emits}. A worse (fp16) head guess can only be rejected, never emitted. Head
precision is therefore purely a speed knob and cannot affect correctness -- which is why moving
the heads to float16 was provably safe.

% ---------------------------------------------------------------
\section{Typical Acceptance and Temperature}
% ---------------------------------------------------------------

Strict greedy accepts a candidate only if it equals the backbone's argmax, so wider or deeper
trees do not help under it -- the extra branches encode non-argmax guesses greedy never takes.
\emph{Typical} acceptance \cite{cai2024} instead accepts a candidate when it is plausible under
the backbone's distribution,
\[
p(\text{cand}) > \min\!\big(\varepsilon,\ \delta\, e^{-H(p)}\big),
\]
where $H(p)$ is the entropy of the backbone's distribution: when the model is unsure the
threshold drops and more candidates pass. This is lossy (it is temperature sampling, not
greedy) but lets deeper and wider trees finally pay off.

Temperature is the knob. Lower $T$ sharpens the distribution and accepts only confident tokens
(coherent, fewer accepts); higher $T$ accepts more but eventually degrades into incoherent text.
This is a constrained optimization -- maximize speed subject to the output staying coherent --
whose optimum sits on the constraint boundary (the highest $T$ that is still coherent). Two
findings: (i) the safe temperature must be tuned against the \emph{worst} prompt, not an average
one ($T{=}0.92$ was coherent on a physics prompt but produced word-salad on a simple list);
(ii) $T{=}0.8$ is robust across prompt types, giving a coherent $\sim$\textbf{1.47x} (1.27--1.67x
by prompt).

% ---------------------------------------------------------------
\section{More Heads Can Be Slower: Optimal-\texorpdfstring{$K$}{K}, Measured}
% ---------------------------------------------------------------

A six-head checkpoint was trained to deepen the fused tree. Under strict greedy it was
\emph{slower} than four heads (1.14x vs.\ 1.24x). The two far-future heads are almost always
wrong, so under greedy they are almost always rejected -- contributing near-zero extra accepted
tokens while still adding tree-verification cost every round. This is the optimal-$K$ tradeoff
(linear, certain cost versus diminishing, conditional acceptance) shown directly on hardware:
past the optimal head count, more heads cost more than they return. The far heads can only earn
their place under typical acceptance, or by spending the budget intelligently -- the subject of
the next section.

% ---------------------------------------------------------------
\section{Calibrated Tree Attention (Original Contribution)}
% ---------------------------------------------------------------

\subsection{The Problem}
The Cartesian tree spends its node budget uniformly: top-2 of every head, every combination
($2+4+8 = 14$ nodes at depths $0,1,2$ for the fused tree, allocated $[2,4,8]$). But the heads
are not equally accurate. Calibration measured the best-guess accuracy of each head as
$21.8\%$, $10.5\%$, $7.3\%$ by depth. The naive tree therefore spends \emph{eight of fourteen
nodes} -- the majority -- on the worst head's guesses, which are almost always rejected.

\subsection{Calibration}
For each tree depth $d$ (which uses head $d{+}1$) and rank $i$, the accuracy
\[
a[d][i] = \Pr\big[\text{head } (d{+}1)\text{'s rank-}i \text{ token equals the true future token}\big]
\]
is estimated on \emph{in-distribution} text -- the model's own greedy continuations of several
prompts -- by applying each head at every position and counting matches against the realized
future token. Accuracy decreases with both depth (further future is harder) and rank.

\subsection{Greedy Node Selection}
A candidate node at depth $d$ reached by rank choices $(i_0,\dots,i_d)$ has path-acceptance
probability equal to the product of accuracies along it,
$P(\text{node}) = \prod_{j=0}^{d} a[j][i_j]$. With a fixed budget $B$ (set to 14, matching the
Cartesian tree), nodes are selected greedily: maintain a frontier of nodes whose parent is
already chosen, repeatedly add the highest-$P$ frontier node and push its children, until $B$
nodes are selected. \textbf{This is Prim's minimum-spanning-tree algorithm} -- grow a tree by
attaching the best-connected frontier node -- with the objective changed from minimizing edge
cost to maximizing cumulative acceptance. Because a child becomes a candidate only after its
parent is chosen, parents precede children, exactly the ordering the decoder requires; the
topology is computed once and stored.

One decoder change supports it: the Cartesian tree's leaves are always its last
$\text{width}^{\text{depth}}$ nodes, but a calibrated tree's leaves can sit at any depth, so
leaves are detected generically (a node that is nobody's parent). The accept rule is unchanged,
so the calibrated greedy path is lossless \emph{by construction} -- the tree only changes which
guesses are proposed, never which token is emitted.

\subsection{Result: the Optimizer Re-discovers Optimal-\texorpdfstring{$K$}{K}}
With 14 nodes, greedy selection allocated $[4, 10, 0]$ across depths -- assigning \textbf{zero
nodes to the weakest head}. Its best guess ($7.3\%$) was too unlikely to be accepted to justify
even one node over a fourth node on the best head. The optimization thus \emph{automatically}
re-derived the optimal-$K$ result of the previous section, with no human intervention. At the
same node budget, greedy and lossless:

\begin{table}[H]
\centering
\caption{Cartesian vs.\ calibrated tree (greedy, lossless), same 14-node budget.}
\begin{tabular}{@{}lccc@{}}
\toprule
Prompt & Cartesian & Calibrated & Acceptance gain \\
\midrule
Relativity & 1.81 tok/round, 47.5 tok/s & 1.91 tok/round, 50.7 tok/s & +6\% \\
Tea (list)  & 1.92 tok/round, 51.0 tok/s & 2.22 tok/round, 58.8 tok/s & +15\% \\
\bottomrule
\end{tabular}
\end{table}

The calibrated tree lifts the lossless speedup from 1.24x to \textbf{1.34--1.56x} with no
temperature, no retraining, and no quality cost -- purely by spending the same budget more
wisely.

% ---------------------------------------------------------------
\section{Results Summary}
% ---------------------------------------------------------------

\begin{table}[H]
\centering
\caption{Full results ladder. Llama 3.1 8B Instruct, RTX A5000, 100 tokens.}
\begin{tabular}{@{}lcc@{}}
\toprule
Configuration & Speedup & Quality \\
\midrule
Naive 8B baseline & 1.00x & -- \\
SpecDecode, 1B draft, $K{=}4$ & 1.17x & lossless (distribution) \\
Medusa greedy, epoch-1 (3-pass) & 0.65x & lossless \\
Medusa 4-head, fused tree, greedy & 1.24x & lossless \\
Medusa 6-head, fused tree, greedy & 1.14x & lossless (optimal-$K$ overshoot) \\
Medusa 4-head, typical, $T{=}1.0$ & 1.76x & incoherent \\
Medusa 4-head, typical, $T{=}0.8$ & $\sim$1.47x & coherent (lossy) \\
\textbf{Medusa 4-head, calibrated tree, greedy} & \textbf{1.34--1.56x} & \textbf{lossless} \\
\bottomrule
\end{tabular}
\end{table}

% ---------------------------------------------------------------
\section{System}
% ---------------------------------------------------------------

The complete system includes a FastAPI server exposing two endpoints:
\texttt{POST /generate} for synchronous generation and \texttt{WS /stream} for
token-by-token streaming. A single-file vanilla JavaScript frontend visualizes
tokens live, colored by acceptance status (blue for accepted draft tokens, dark red
for backbone corrections). The WebSocket sends each token as
\texttt{\{text, accepted\}} JSON so the frontend can color without any additional
state.

% ---------------------------------------------------------------
\section{What I Learned}
% ---------------------------------------------------------------

A few things that were not obvious before building this:

\textbf{Model size does not determine draft cost.} The 1B draft costs 40\% of the
8B target per token, not the expected 12.5\%. Single-token forward passes are
dominated by fixed overhead (kernel launch, memory bandwidth for weights), not by
multiply-accumulate FLOPs. This is why the cost model uses $c = 0.40$ instead of
$c = 0.125$.

\textbf{Draft-target dialect matching matters as much as size.} Switching from the
base 1B to the instruct 1B raised K=4 acceptance from 2.45 to 3.05 tokens/round --
a bigger gain than any hyperparameter tuning. Two models with different training
objectives operate in different output distributions even at the same positions.

\textbf{The DynamicCache mutation bug.} Transformers' \texttt{DynamicCache} is
modified in place during every forward pass. Passing the same object to two
sequential calls contaminates the second call's computation. The fix -- serializing
and deserializing through the legacy tuple format to create a fresh object --
is not documented anywhere obvious. Found it by debugging why the cache version
and the non-cache version diverged after the first round.

\textbf{Float16 NaN in Medusa training.} Training Medusa heads in float16 produces
NaN from epoch 0. The softmax over 128,256 tokens computes $e^{x}$ for logit $x$;
at early training, logits for some tokens exceed 11 and $e^{11} \approx 60,000$,
which overflows float16's maximum of 65,504. The fix is mixed precision: backbone
stays float16 (no gradients), heads train in float32.

\textbf{Correctness testing under fp16 must not be bit-exact.} A near-tie in the
logits resolves to different argmaxes depending on the reduction order of the
forward pass (prefill vs.\ decode). A decoder can be a perfectly valid greedy decode
yet differ from a reference implementation token-for-token. The right correctness
test is self-consistency (is every emitted token argmax-or-tied?), not equality to a
reference -- a subtle point that took a full debugging session to learn.

\textbf{An optimizer can re-derive a hand-found result.} The optimal-$K$ tradeoff
was first observed by hand: a sixth head made greedy slower. The calibrated tree's
greedy node-selection then re-derived it automatically, assigning zero nodes to the
weakest head. Framing the tree construction as ``maximize expected accepted tokens
subject to a node budget'' turned a heuristic into an explicit optimization whose
solution \emph{contained} the insight -- the most satisfying moment of the project,
and the clearest place where the operations-research framing paid off.

% ---------------------------------------------------------------
\section{What I Would Do Next}
% ---------------------------------------------------------------

\begin{itemize}
    \item \textbf{Adaptive draft length.} Treat the head count / tree depth as a control
          variable adjusted online from a sliding window of recent acceptance, increasing it
          on predictable text and shrinking it on hard text. This is the same throughput
          optimization that the calibrated tree solves \emph{offline}, made \emph{online} --
          a natural operations-research extension.
    \item \textbf{Online / per-context calibration.} The calibration here is global. Accuracy
          varies by domain (code vs.\ prose); recalibrating the tree topology per context, or
          maintaining a few topologies and selecting by prompt type, should add acceptance.
    \item \textbf{Batched speculative decoding.} Verifying trees for many sequences in one
          backbone pass, to raise throughput under concurrent load -- the production setting.
    \item \textbf{MT-Bench quality evaluation.} The greedy paths are lossless by construction
          and verified by self-consistency; the typical path is lossy, and a formal MT-Bench
          comparison would quantify its quality cost against the speed gain.
    \item \textbf{Public deployment.} Host the live visualization behind a small GPU so the
          colour-coded token stream is reachable from a URL (benchmark numbers remain those
          measured on the A5000).
\end{itemize}

% ---------------------------------------------------------------
\section{References}
% ---------------------------------------------------------------

\begin{thebibliography}{9}

\bibitem{leviathan2023}
Yaniv Leviathan, Matan Kalman, and Yossi Matias.
\textit{Fast Inference from Transformers via Speculative Decoding.}
ICML 2023. \href{https://arxiv.org/abs/2211.17192}{arXiv:2211.17192}

\bibitem{cai2024}
Tianle Cai, Yuhong Li, Zhengyang Geng, Hongwu Peng, Jason D. Lee,
Deming Chen, and Tri Dao.
\textit{Medusa: Simple LLM Inference Acceleration Framework with
Multiple Decoding Heads.}
ICML 2024. \href{https://arxiv.org/abs/2401.10774}{arXiv:2401.10774}

\bibitem{dao2022}
Tri Dao, Daniel Y. Fu, Stefano Ermon, Atri Rudra, and Christopher Re.
\textit{FlashAttention: Fast and Memory-Efficient Exact Attention
with IO-Awareness.}
NeurIPS 2022. \href{https://arxiv.org/abs/2205.14135}{arXiv:2205.14135}

\end{thebibliography}

\end{document}
